\documentclass[10pt,a4paper]{article}

\usepackage[a4paper,top=13mm,bottom=14mm,left=14mm,right=14mm]{geometry}
\usepackage{fontspec}
\usepackage{xcolor}
\usepackage{tabularx}
\usepackage{array}
\usepackage[most]{tcolorbox}
\usepackage{fancyhdr}
\usepackage{enumitem}
\usepackage{amssymb}
\usepackage{microtype}

\IfFontExistsTF{Arial}{
  \setmainfont{Arial}
  \setsansfont{Arial}
}{
  \setmainfont{TeX Gyre Heros}
  \setsansfont{TeX Gyre Heros}
}

\definecolor{AIEPRed}{HTML}{D62027}
\definecolor{AIEPNavy}{HTML}{082B5C}
\definecolor{AIEPInk}{HTML}{182B3A}
\definecolor{AIEPMuted}{HTML}{5F6B7A}
\definecolor{AIEPLine}{HTML}{C9D2DE}
\definecolor{AIEPSoft}{HTML}{F5F2EC}
\definecolor{AIEPBlueSoft}{HTML}{E9EEF4}
\definecolor{AIEPCyan}{HTML}{35B7C6}
\definecolor{AIEPGold}{HTML}{E0BC5A}
\definecolor{AIEPGreen}{HTML}{2E8B57}

\pagestyle{fancy}
\fancyhf{}
\renewcommand{\headrulewidth}{0pt}
\fancyfoot[L]{\footnotesize\textcolor{AIEPMuted}{GeoGreen Escolar · Mentoría 4}}
\fancyfoot[R]{\footnotesize\textcolor{AIEPMuted}{Ficha de trabajo por equipo · \thepage/4}}

\setlength{\parindent}{0pt}
\setlength{\parskip}{3pt}
\renewcommand{\arraystretch}{1.18}
\setlist[itemize]{leftmargin=*,topsep=1pt,itemsep=1pt}
\newcolumntype{Y}{>{\raggedright\arraybackslash}X}
\newcolumntype{C}[1]{>{\centering\arraybackslash}p{#1}}

\newtcolorbox{fieldbox}[2][]{
  colback=white,
  colframe=AIEPLine,
  boxrule=0.55pt,
  arc=1mm,
  left=2.5mm,right=2.5mm,top=2mm,bottom=2mm,
  before skip=2pt,after skip=2pt,
  title={\scriptsize\bfseries\MakeUppercase{#2}},
  coltitle=AIEPNavy,
  fonttitle=\sffamily,
  attach title to upper={\par\vspace{0.8mm}},
  #1
}

\newtcolorbox{accentfield}[2][]{
  colback=AIEPSoft,
  colframe=AIEPRed,
  boxrule=0.7pt,
  arc=1mm,
  left=2.5mm,right=2.5mm,top=2mm,bottom=2mm,
  before skip=2pt,after skip=2pt,
  title={\scriptsize\bfseries\MakeUppercase{#2}},
  coltitle=AIEPRed,
  fonttitle=\sffamily,
  attach title to upper={\par\vspace{0.8mm}},
  #1
}

\newcommand{\writearea}[1]{\vspace{#1}\hrule height 0.35pt\vspace{4mm}\hrule height 0.35pt}
\newcommand{\checkline}[1]{$\square$\hspace{1.5mm}#1}
\newcommand{\pagehead}[2]{%
  {\small\bfseries\textcolor{AIEPRed}{GEOGREEN ESCOLAR · MENTORÍA 4}}\par
  \vspace{2mm}
  {\LARGE\bfseries\color{AIEPNavy}#1\par}
  {\normalsize\color{AIEPMuted}#2\par}
  \vspace{4mm}
}

\begin{document}

\pagehead{Recuperar la historia del proyecto}{Ficha de trabajo por equipo · Conservar como base de la presentación}

\begin{tabularx}{\textwidth}{@{}>{\bfseries\color{AIEPNavy}}p{25mm}Y >{\bfseries\color{AIEPNavy}}p{29mm}p{38mm}@{}}
Equipo & \rule{57mm}{0.35pt} & Proyecto & \rule{36mm}{0.35pt} \\
Integrantes & \multicolumn{3}{l}{\rule{145mm}{0.35pt}} \\
\end{tabularx}

\vspace{3mm}
\begin{tcolorbox}[colback=AIEPNavy,colframe=AIEPNavy,coltext=white,boxrule=0pt,arc=1mm,left=3mm,right=3mm,top=2mm,bottom=2mm]
\textbf{Productos de entrada:}\hspace{4mm}
\checkline{Mentoría 1}\hspace{5mm}
\checkline{Mentoría 2}\hspace{5mm}
\checkline{Mentoría 3}
\end{tcolorbox}

{\large\bfseries\color{AIEPNavy}1. Inventario narrativo}\par
{\small Escriban una frase breve por anclaje. Si falta respaldo, registren ``falta comprobar''.}

\begin{tabularx}{\textwidth}{@{}Y@{\hspace{4mm}}Y@{}}
\begin{fieldbox}[height=42mm]{Problema y contexto}
¿Qué ocurre, dónde ocurre y quiénes se relacionan con la situación?
\writearea{10mm}
\end{fieldbox}
&
\begin{fieldbox}[height=42mm]{Solución propuesta}
¿Qué propone el equipo y de qué manera responde al problema?
\writearea{10mm}
\end{fieldbox}
\\
\begin{fieldbox}[height=42mm]{Evidencia de avance}
¿Qué construimos, probamos, representamos o aprendimos?
\writearea{10mm}
\end{fieldbox}
&
\begin{fieldbox}[height=42mm]{Aprendizaje o aporte esperado}
¿Qué comprendimos mejor y qué mejora busca favorecer la propuesta?
\writearea{10mm}
\end{fieldbox}
\end{tabularx}

\begin{accentfield}[height=42mm]{Idea principal que queremos dejar}
Si la audiencia recordara una sola idea, debería recordar que:
\writearea{12mm}
\end{accentfield}

\begin{tcolorbox}[colback=AIEPBlueSoft,colframe=AIEPLine,boxrule=0.5pt,arc=1mm,left=3mm,right=3mm,top=2mm,bottom=2mm]
\textbf{Comprobación:} \checkline{Las cuatro frases se conectan.}\quad
\checkline{Podemos explicar de dónde salió cada una.}\quad
\checkline{No presentamos expectativas como resultados.}
\end{tcolorbox}

\newpage
\pagehead{Construir el guion del pitch}{Una idea principal por tramo · Frases breves · Información que el equipo pueda explicar}

{\large\bfseries\color{AIEPNavy}2. Orden de la historia}\par
{\small Completen con palabras clave. El guion orienta la explicación; no es un texto para leer de manera continua.}

\small
\begin{tabularx}{\textwidth}{|C{13mm}|p{34mm}|p{52mm}|Y|}
\hline
\textbf{N.º} & \textbf{Tramo} & \textbf{Pregunta que debe responder} & \textbf{Idea principal del equipo} \\
\hline
\rule{0pt}{18mm}1 & Problema & ¿Qué ocurre y en qué contexto? & \\
\hline
\rule{0pt}{18mm}2 & Personas & ¿A quiénes involucra o afecta? & \\
\hline
\rule{0pt}{18mm}3 & Solución & ¿Qué proponemos y por qué? & \\
\hline
\rule{0pt}{18mm}4 & Funcionamiento & ¿Qué hace la propuesta durante su uso? & \\
\hline
\rule{0pt}{18mm}5 & Evidencia & ¿Qué podemos mostrar y qué demuestra? & \\
\hline
\rule{0pt}{18mm}6 & Aporte y cierre & ¿Qué busca mejorar y qué idea debe recordarse? & \\
\hline
\end{tabularx}
\normalsize

\vspace{4mm}
\begin{accentfield}[height=35mm]{Apertura del pitch}
En \rule{36mm}{0.35pt} observamos que \rule{92mm}{0.35pt}.\par
Esta situación involucra o afecta a \rule{85mm}{0.35pt}.
\end{accentfield}

\begin{accentfield}[height=35mm]{Cierre del pitch}
Podemos respaldar nuestro avance mostrando \rule{80mm}{0.35pt}.\par
Con esta propuesta esperamos aportar a \rule{83mm}{0.35pt}.
\end{accentfield}

\begin{tcolorbox}[colback=AIEPNavy,colframe=AIEPNavy,coltext=white,boxrule=0pt,arc=1mm,left=3mm,right=3mm,top=2.5mm,bottom=2.5mm]
\textbf{Prueba de comprensión}\par
Sin mirar esta ficha, otra persona debe poder responder: ¿cuál es el problema?, ¿qué propone el equipo?, ¿cómo funciona? y ¿qué evidencia existe? Si necesita una explicación adicional, revisen el orden o la precisión del tramo.
\end{tcolorbox}

\newpage
\pagehead{Preparar la presentación compartida}{Soporte visual · Seis voces · Acciones · Transiciones · Tiempos}

{\large\bfseries\color{AIEPNavy}3. Soporte visual con propósito}\par
\begin{tabularx}{\textwidth}{@{}Y@{\hspace{4mm}}Y@{}}
\begin{fieldbox}[height=37mm]{Recurso elegido}
¿Qué mostraremos: imagen, diagrama, video, maqueta, simulación, prototipo u otro recurso?
\writearea{6mm}
\end{fieldbox}
&
\begin{fieldbox}[height=37mm]{Función comunicativa}
Marquen una: \checkline{situar}\quad\checkline{explicar}\quad\checkline{demostrar}
\writearea{6mm}
\end{fieldbox}
\end{tabularx}

\begin{fieldbox}[height=32mm]{Qué debe observar y comprender la audiencia}
No basta con mostrar el recurso: indiquen qué evidencia contiene y por qué importa.
\writearea{5mm}
\end{fieldbox}

{\large\bfseries\color{AIEPNavy}4. Pauta de seis intervenciones}\par
{\small Cada integrante registra una idea sustantiva, una acción y una frase que entregue el hilo a quien continúa.}

\scriptsize
\begin{tabularx}{\textwidth}{|C{8mm}|p{25mm}|p{36mm}|p{30mm}|p{17mm}|Y|}
\hline
\textbf{N.º} & \textbf{Rol e integrante} & \textbf{Idea principal} & \textbf{Soporte o acción} & \textbf{Tiempo} & \textbf{Transición de salida} \\
\hline
\rule{0pt}{18mm}1 & Coordinación: & & & & \\
\hline
\rule{0pt}{18mm}2 & Investigación: & & & & \\
\hline
\rule{0pt}{18mm}3 & Diseño: & & & & \\
\hline
\rule{0pt}{18mm}4 & Tecnología: & & & & \\
\hline
\rule{0pt}{18mm}5 & Pruebas y evidencia: & & & & \\
\hline
\rule{0pt}{18mm}6 & Comunicación: & & & & \\
\hline
\end{tabularx}
\normalsize

\begin{tcolorbox}[colback=AIEPBlueSoft,colframe=AIEPLine,boxrule=0.5pt,arc=1mm,left=3mm,right=3mm,top=2mm,bottom=2mm]
\textbf{Antes de ensayar:} \checkline{Las seis personas aportan contenido.}\quad
\checkline{El soporte tiene responsable.}\quad
\checkline{Las transiciones conectan ideas.}\quad
\checkline{Todos conocen la secuencia completa.}
\end{tcolorbox}

\newpage
\pagehead{Ensayar, retroalimentar y corregir}{Ejecutar una versión · Observar acciones · Cambiar una parte · Volver a probar}

{\large\bfseries\color{AIEPNavy}5. Pauta común de observación}\par
\begin{tabularx}{\textwidth}{@{}Y@{\hspace{3mm}}Y@{}}
\begin{fieldbox}[height=31mm]{Se entiende}
¿Problema, contexto e idea principal usan lenguaje claro?\par
\checkline{Sí}\quad\checkline{Aún no}
\end{fieldbox}
&
\begin{fieldbox}[height=31mm]{Se conecta}
¿Problema, solución y aporte mantienen una relación coherente?\par
\checkline{Sí}\quad\checkline{Aún no}
\end{fieldbox}
\\
\begin{fieldbox}[height=31mm]{Se demuestra}
¿La evidencia se muestra, se explica y respalda lo afirmado?\par
\checkline{Sí}\quad\checkline{Aún no}
\end{fieldbox}
&
\begin{fieldbox}[height=31mm]{Se coordina}
¿Roles, soporte, transiciones y tiempos funcionan como equipo?\par
\checkline{Sí}\quad\checkline{Aún no}
\end{fieldbox}
\end{tabularx}

{\large\bfseries\color{AIEPNavy}6. Retroalimentación accionable}\par
\begin{tabularx}{\textwidth}{@{}Y@{\hspace{3mm}}Y@{\hspace{3mm}}Y@{}}
\begin{fieldbox}[height=41mm]{Observación}
¿Qué ocurrió de manera concreta?
\writearea{8mm}
\end{fieldbox}
&
\begin{fieldbox}[height=41mm]{Efecto}
¿Qué facilitó o dificultó comprender?
\writearea{8mm}
\end{fieldbox}
&
\begin{fieldbox}[height=41mm]{Ajuste}
¿Qué acción específica volveremos a probar?
\writearea{8mm}
\end{fieldbox}
\end{tabularx}

\begin{accentfield}[height=38mm]{Corrección comprobada}
Antes, \rule{55mm}{0.35pt}. Cambiamos \rule{68mm}{0.35pt}.\par
Al reensayar, observamos que \rule{99mm}{0.35pt}.
\end{accentfield}

{\large\bfseries\color{AIEPNavy}7. Lista de salida obligatoria}\par
\begin{tcolorbox}[colback=AIEPNavy,colframe=AIEPNavy,coltext=white,boxrule=0pt,arc=1mm,left=3mm,right=3mm,top=2.5mm,bottom=2.5mm]
\checkline{Guion del pitch actualizado}\hspace{5mm}
\checkline{Soporte visual identificado}\par
\checkline{Seis roles, transiciones y tiempos}\hspace{5mm}
\checkline{Registro de mejora incorporado}
\end{tcolorbox}

\begin{fieldbox}[height=32mm]{Versión que continuará preparando el equipo}
La idea que queremos que la audiencia recuerde es:
\writearea{5mm}
\end{fieldbox}

\end{document}
